% Generated draft: requires mathematical and source review.
\documentclass[11pt]{article}
\usepackage[margin=25mm]{geometry}
\usepackage{fontspec}
\setmainfont{Latin Modern Roman}
\usepackage{amsmath,amssymb,mathtools,unicode-math}
\setmathfont{Latin Modern Math}
\usepackage{enumitem,needspace,xurl,hyperref}
\setlength{\parindent}{0pt}
\setlength{\parskip}{5pt}
\setlength{\emergencystretch}{3em}
\newcommand{\N}{\mathbb{N}}
\newcommand{\Z}{\mathbb{Z}}
\newcommand{\Q}{\mathbb{Q}}
\newcommand{\R}{\mathbb{R}}
\newcommand{\C}{\mathbb{C}}
\newcommand{\F}{\mathbb{F}}
\newcommand{\A}{\mathbb{A}}
\newcommand{\PP}{\mathbb{P}}
\newcommand{\E}{\mathbb{E}}
\newcommand{\eps}{\varepsilon}
\newcommand{\dd}{\,\mathrm{d}}
\begin{document}
\section*{UnsolvedMath ID: 800488}
\textbf{Source problem code:} catalogue entry
\subsection{Definitions and mathematical statement}
For the three-dimensional benchmark, let \(M\) be a closed connected smooth \(3\)-manifold and let \(\alpha\) be a contact form, meaning that \(\alpha\wedge\dd\alpha\) is nowhere zero. Its Reeb vector field \(R_\alpha\) is determined by
\[
\alpha(R_\alpha)=1,\qquad \dd\alpha(R_\alpha,\cdot)=0.
\]
A periodic orbit is the image of a nonconstant periodic trajectory. For purposes of stating the benchmark, use the convention
\[
P_\alpha(T)=\#\{\text{geometrically distinct periodic orbits of }R_\alpha\text{ with least period at most }T\},\qquad T>0.
\]
Thus iterates of the same orbit are not counted separately. The supplied excerpt does not define \(P(T)\); this counting convention requires source confirmation. If a count is infinite, the displayed growth expression is interpreted using \(\log(+\infty)=+\infty\).

With that convention, the contextual conjecture reproduced in the source is: every Reeb flow on a closed connected \(3\)-manifold either has exactly two geometrically distinct periodic orbits or satisfies
\[
\liminf_{T\to+\infty}\frac{\log P_\alpha(T)}{\log T}\geq 2.
\]
The source explicitly states that proving this three-dimensional conjecture is \emph{not} the proposed problem.

The proposed problem is to investigate whether the same phenomenon holds in every dimension. In the Reeb setting, the standard higher-dimensional objects would be contact forms \(\alpha\) on closed connected manifolds \(M^{2n+1}\), with
\[
\alpha\wedge(\dd\alpha)^n\neq0,
\]
and their Reeb flows, with periodic orbits counted as above. However, the supplied record does not specify the higher-dimensional dichotomy: it does not state whether the exceptional orbit count remains exactly two, whether that count depends on dimension, whether the growth exponent remains two, or what additional hypotheses are intended. Consequently, the source supports an investigation of a higher-dimensional analogue, but does not determine a unique quantified conjecture in higher dimensions.

There is a further scope qualification. Before posing the higher-dimensional question, the source also discusses general low-dimensional symplectic dynamics: symplectic diffeomorphisms of compact symplectic surfaces, with or without boundary, and Hamiltonian flows on compact regular \(3\)-dimensional energy levels without boundary. It is unclear whether the proposed extension encompasses higher-dimensional analogues of both these classes as well as Reeb flows. No precise higher-dimensional statement for these additional classes is supplied.

The following are supplied background assessments, not independently verified claims of this draft. The source says that the three-dimensional conjecture is also cited in [DanCG], and that [Albach] confirms it for geodesic flows of reversible Finsler metrics on \(S^2\). It compares this with the earlier estimate attributed to Hingston [Hingston1993],
\[
\liminf_{T\to+\infty}P(T)\frac{\log T}{T}>0.
\]
It states that the number of closed geodesics in this sphere case is always infinite, so that the exceptional finite-orbit alternative does not occur. It further reports that the conjecture has long been known for geodesic flows on closed connected orientable reversible Finsler surfaces of genus at least one, and concludes that every closed reversible Finsler surface, orientable or not, has what it calls quadratic growth of closed geodesics. The displayed logarithmic growth inequality is the explicit benchmark here; the phrase quadratic growth alone does not specify a uniform lower bound of the form \(cT^2\).

The source also states that geodesic flows on \(S^2\) lift to Reeb flows on the standard contact \(3\)-sphere, and that [Albach] treats a significantly larger class of such Reeb flows. Finally, it attributes to [Albach] the same dichotomy for every symplectic diffeomorphism of \(S^2\), using an argument based on transverse foliations of Le Calvez and Tal [LeCalvez-Tal]. These contextual assertions do not resolve the formulation of the proposed higher-dimensional problem.
\subsection{Short English statement}
Investigate a higher-dimensional analogue of the periodic-orbit dichotomy discussed for Reeb flows in dimension three: exactly two periodic orbits, or logarithmic orbit-growth exponent at least two. The source explicitly excludes proving the three-dimensional conjecture as the proposed task. It leaves the higher-dimensional exceptional count, growth threshold, hypotheses, and dynamical scope unspecified. The registry labels this problem open; this is a statement-only draft requiring human mathematical review.
\subsection{Sources}
\begin{itemize}
\item Supplied local UnsolvedMath registry record: permanent ID \(800488\), source code \texttt{catalogue entry}, collection \emph{AMR Open Problem Lists}, upstream ID \(800488\), upstream revision \texttt{d371ee917ac9647fbb06cd6fc529deec6f3dda82}. The record supplies the status \emph{open}, proposed year \(2026\), and proposal attribution to Albach, Barraud et al.; these are registry metadata, not independently verified mathematical assessments.
\item Supplied source list: Albach, Barraud et al., \emph{Open Problems in Billiards and Quantitative Symplectic Geometry} (2026), \url{https://arxiv.org/abs/2602.12896}. The specific source item is attributed to Bernhard Albach, Umberto Hryniewicz, and Alexandru Oancea, Problem 1; the same names are listed as contributors. The supplied provenance records access on 2026-10-16 and extraction by \texttt{arxiv\_latex}. This access information is supplied provenance, not a claim of independent source verification.
\item The supplied excerpt cites [DanCG] for another occurrence of Conjecture 1 and [Albach] for the reported results on reversible Finsler geodesic flows, a larger class of Reeb flows, and symplectic diffeomorphisms of \(S^2\). No full bibliographic entries or separate links for these labels are supplied.
\item The supplied excerpt attributes the earlier growth estimate to Hingston [Hingston1993] and the transverse-foliation input to Le Calvez and Tal [LeCalvez-Tal]. No full bibliographic entries or separate links for these labels are supplied.
\end{itemize}
\end{document}
